\section{Extraction of the prime statistic}
\label{ext:section}

Fix $d\geq1$ and $\eps>0$.  For each $i\leq d$, let
$\mathcal I_i$ be an interval of primes with lower endpoint at least
$x^\eps$, and suppose that the product of the upper endpoints is
at most $x^{1-\eps}$.  Define
\begin{equation}\label{ext:statistic}
 f_x(u)=\sum_{\substack{\ell_i\in\mathcal I_i\\
                         \ell_1,\ldots,\ell_d\text{ distinct}}}
                   \1_{\ell_1\cdots\ell_d\mid u},\qquad
 C_x=\sum_{\substack{\ell_i\in\mathcal I_i\\
                         \ell_1,\ldots,\ell_d\text{ distinct}}}
                   \frac1{\ell_1\cdots\ell_d}.
\end{equation}
Lemma~\ref{lem:prime-estimates} gives $C_x=O_{d,\eps}(1)$.
On any fixed range $u\asymp x$, the same bound holds for $f_x(u)$:
there are only $O_\eps(1)$ prime divisors of $u$ exceeding
$x^\eps$.  All cutoffs below are supported in a fixed compact
subset of $(0,\infty)$.

\begin{theorem}[Interior prime statistic]\label{thm:interior}
For every nonnegative $\Phi\in C_c^\infty((0,\infty))$,
\begin{equation}\label{eq:interior}
 \sum_{p\text{ prime}}\Phi((p-1)/x)
                   \bigl(f_x(p-1)-C_x\bigr)=o(x/\log x).
\end{equation}
The limit holds through all real $x\to\infty$ for the slot
intervals just described.
\end{theorem}

Throughout the proof write $L=\log x$, $W=\exp(L^{.24})$, and
$V(W)=\prod_{p\leq W}(1-1/p)$.  Fix $q,q_0\in(0,1)$, for
example $q=q_0=1/2$.  A \emph{marking array} will mean $K$ disjoint
bands from Theorem~\ref{thm:type-II}, with their weight
\begin{equation}\label{ext:array-weight}
 \mathcal W(u)=\prod_{i=1}^K
          q^{\omega_i(u)-1}\frac{\omega_i(u)}{V_i},\qquad
 V_i=\sum_{p\in\mathcal P_i}\frac1p.
\end{equation}
Each band consists of all primes between $\exp(L^{a_i})$ and
$\exp(2L^{a_i})$, with $.1<a_i<.2$.  The number $K$ and the
exponents are fixed in each limit.  Since $V_i=\log2+o(1)$ and
$\sup_{n\geq0}nq^{n-1}<\infty$, these weights are bounded by a
constant depending only on $K,q$, for large $x$.

Presieving $u+1$ leaves both primes and composites in the average of
$\mathcal W(u)(f_x(u)-C_x)$.  We will show that the composite
contribution is negligible in two stages.  The one-sided estimate
first replaces its prime factors by rough variables; additional
divisor marks then allow the two-sided estimate to bound the resulting
centered sum.  This isolates a prime average carrying $\mathcal W$.
Finally, a mean-square bound for averages of disjoint marking arrays
removes that weight.

\subsection{Progressions and the independent divisor model}

For finitely many disjoint arrays, define the independent divisor
model by replacing every $\1_{p\mid u}$ at a band prime with an
independent Bernoulli variable of mean $1/p$.  For a function $H$
of these indicators write $M_x(H)$ for its expectation in this
model.  In particular put $m_{\mathcal W}=M_x(\mathcal W)$.
There are constants $0<c_K<C_K<\infty$ such that
\begin{equation}\label{ext:model-bounds}
 c_K\leq m_{\mathcal W}\leq C_K,
 \qquad M_x(\mathcal W^2)\leq C_K.
\end{equation}
For the lower bound, in each band the event of exactly one hit has
probability
$\prod_p(1-1/p)\sum_p1/(p-1)$, bounded below since the reciprocal
sum tends to $\log2$ and the largest $1/p$ tends to zero.
The weights have a fixed positive value on the event of one hit
per band.  Independence between bands proves the lower bound;
boundedness proves the other two assertions.  These constants
can be chosen independently of the particular fixed exponents.

\begin{lemma}[Summed progression remainder]\label{lem:progressions}
Choose $0<c_0<\eps/2$.  Let $H$ be a weight from a fixed finite
collection of disjoint arrays, a product of two such weights
(allowing a repeated array), a constant, or a bounded linear
combination of these functions.  Then, for every fixed $A>0$,
\begin{align}\label{ext:progression}
 \sum_{r\leq x^{c_0}}\left|
  \sum_{\substack{u\geq1\\r\mid u+1}}
             \Phi(u/x)f_x(u)H(u)
  -\frac{x}{r}\left(\int_0^\infty\Phi(t)\,\dd t\right)
                     C_x M_x(H)\right|
  &\ll_A xL^{-A}.
\end{align}
The same assertion holds with $f_x,C_x$ replaced by $1,1$.
The constants may depend on the fixed arrays and on the bounded
coefficients in the linear combination.
\end{lemma}

\begin{proof}
We give the expansion including the repeated-array case needed
later.  It suffices to treat one product of weights.  In any band
$g$ which occurs, let $a_g\in\{1,2\}$ be its multiplicity in the
product.  Its factor is
\[
 V_g^{-a_g}q^{a_g(\omega_g(u)-1)}\omega_g(u)^{a_g}.
\]
Expand $\omega_g^{a_g}$ as the sum over ordered $a_g$-tuples of
divisor primes, permitting coincidences.  For a chosen tuple let
$S_g$ be its set of distinct primes.  On the event $S_g\mid u$,
where this notation means that every prime in $S_g$ divides $u$,
the contribution is
\begin{equation}\label{ext:represented-union}
 V_g^{-a_g}q^{a_g(|S_g|-1)}\1_{S_g\mid u}
 \prod_{p\in\mathcal P_g\setminus S_g}
        \bigl(1-(1-q^{a_g})\1_{p\mid u}\bigr).
\end{equation}
Thus repetitions change the damping to $q^2$ on the unrepresented
primes.  For instance, the reciprocal mass of the represented
unions when $a_g=2$ is exactly
\[
 V_g^{-2}\left\{V_g+q^2\left(V_g^2-
                                  \sum_{p\in\mathcal P_g}p^{-2}\right)\right\},
\]
and is bounded.  For $a_g=1$ that mass is one.
Multiplying over the fixed number of bands shows that all
represented unions have total reciprocal mass $O(1)$, with their
nonnegative coefficients and their multiplicities included.

Expand the remaining product in \eqref{ext:represented-union}
through degree $h_0=\lfloor L^{.01}\rfloor$ in the unrepresented
primes.  Bonferroni inequalities, valid for factors $1-t_p$ with
$0\leq t_p\leq1$, bound the error in absolute value by the
nonnegative term of degree $h_0+1$.  If $C_*$ bounds the sum of
the reciprocal masses of all bands involved, the total reciprocal
coefficient mass at degree $t$ is at most a fixed constant times
$C_*^t/t!$.  Hence the truncated expansion has reciprocal
coefficient mass $O(1)$ and its error envelope has reciprocal
mass at most
\begin{equation}\label{ext:bonferroni-tail}
 O\left(\frac{C_*^{h_0+1}}{(h_0+1)!}\right)=O_A(L^{-A})
 \quad\text{for every fixed }A.
\end{equation}
Every divisor in the truncation or the error envelope satisfies
\begin{equation}\label{ext:model-divisor-size}
 d'\leq\exp\bigl(O_K((h_0+1)L^{.2})\bigr)
       \leq\exp(O_K(L^{.22}))=x^{o(1)}.
\end{equation}
It follows also that the ordinary absolute coefficient mass is
$x^{o(1)}$: multiply its reciprocal mass by the largest divisor.
All these facts hold as well for the constant function.

Fix an ordered long-prime tuple, put $\ell=\prod_i\ell_i$, and
write $u=\ell y$.  Every band prime is coprime to $\ell$.
Moreover, every $\ell_i>x^{c_0}$, so $(\ell,r)=1$ for
$r\leq x^{c_0}$.  The congruence $r\mid\ell y+1$ therefore fixes
one reduced residue class of $y$ modulo $r$.  For a divisor $d'$
of band primes the additional condition $d'\mid y$ is impossible
if $(d',r)>1$.  Otherwise the Chinese remainder theorem and smooth
summation in one progression give
\begin{equation}\label{ext:crt-count}
 \sum_{\substack{y\geq1\\\ell y\equiv-1\ (r)\\d'\mid y}}
        \Phi(\ell y/x)
 =\frac{x}{\ell r d'}\int_0^\infty\Phi(t)\,\dd t+O_\Phi(1).
\end{equation}
The error follows, for example, by comparing each sampling interval
with its integral and using the bounded total variation of $\Phi$.
It is uniform in $\ell,r,d'$.

Inserting the expansion into \eqref{ext:crt-count} recovers the
model in which the band primes dividing $r$ are forced absent.
The same calculation for the positive Bonferroni envelope bounds
the truncation error by \eqref{ext:bonferroni-tail} times
$x/(\ell r)$, plus $x^{o(1)}$ in rounding errors.
There are at most $d!\,x^{1-\eps}$ ordered long-prime tuples,
since a squarefree product determines at most $d!$ orders.
Consequently the sum of all rounding errors over tuples and
$r\leq x^{c_0}$ is
\begin{equation}\label{ext:rounding-sum}
 O(x^{1-\eps+c_0+o(1)})=O(x^{1-\eps/3}),
\end{equation}
after increasing the threshold for $x$.  When $f_x=1$, use
$\ell=1$ and the stronger bound $x^{c_0+o(1)}$ instead.
The principal Bonferroni errors sum to at most
$O(x C_x L C_*^{h_0+1}/(h_0+1)!)$ and save every fixed log power.

Finally, coupling the forced-absence model to the unconditional
one and using the boundedness of $H$ bounds their mean difference
by
\[
 O\left(\sum_{\substack{p\mid r\\p\text{ in the arrays}}}\frac1p\right)
 \ll L\exp(-L^{.1}).
\]
Its contribution after summing $1/\ell$ and $1/r$ is smaller than
every $xL^{-A}$.  This proves \eqref{ext:progression}; bounded
linear combinations follow by the triangle inequality.
\end{proof}

\subsection{Presieving and replacement of composite prime slots}

Fix a large even integer $h$, and set
\begin{equation}\label{ext:sieve-parameters}
 \gamma=\frac{c_0}{4h+3},\qquad z=x^\gamma.
\end{equation}
Apply the block sieve, Lemma~\ref{lem:block-sieve}, to the bad
conditions $r'\mid u+1$ at primes $r'\leq z$, with density
$g(r')=1/r'$.  Use separately the two nonnegative weights
$\Phi(u/x)\mathcal W(u)f_x(u)$ and
$\Phi(u/x)\mathcal W(u)$.  The density is at most $1/2$, and
its reciprocal mass in every block $(b,b^2]$ is bounded by an
absolute constant by Lemma~\ref{lem:prime-estimates}.  The sieve
remainder is bounded by the sum in Lemma~\ref{lem:progressions},
because
\begin{equation}\label{ext:sieve-level}
 z^{4h+2}=x^{c_0(4h+2)/(4h+3)}<x^{c_0}.
\end{equation}
Subtract the second sieve formula times $C_x$ from the first.
Their main terms cancel, while each relative sieve error is
$O(e^{-h})$.  Since the main terms both contain
$m_{\mathcal W}$ and
$\prod_{r'\leq z}(1-1/r')\asymp1/(\gamma L)$, we obtain
\begin{equation}\label{ext:presift}
 \limsup_{x\to\infty}\frac{L}{xm_{\mathcal W}}
 \left|\sum_{\substack{u\geq1\\P^-(u+1)>z}}
    \Phi(u/x)\mathcal W(u)(f_x(u)-C_x)\right|
 \ll\frac{e^{-h}}{\gamma}.
\end{equation}
The implied constant depends on the slot and cutoff data but is
independent of $K,h$ and of the fixed array.  The progression
remainders may have constants depending on these choices; their
normalized limits are zero.  In particular, no reciprocal of a
rare-mark mean occurs in the surviving error in
\eqref{ext:presift}.

For integers in this sum that are composite, write their ordered
prime factorization as $u+1=p_1\cdots p_j$, with $p_i>z$.
For large $x$, $2\leq j\leq2/\gamma$.  Integers divisible by
$p^2$ for a prime $p>z$ number
\[
 \ll\sum_{p>z}\frac{x}{p^2}\ll x/z.
\]
The weights and the possible ordered factorization counts are
bounded for fixed $\gamma,d,K$.  We may therefore discard these
integers at error $o(x/L)$, count the others with weight $1/j!$,
and restore repeated factors when convenient at the same cost.
For each fixed $j$ we will prove
\begin{equation}\label{ext:composite-vanishing}
 \sum_{p_1,\ldots,p_j>z}
 \Phi((p_1\cdots p_j-1)/x)
 \mathcal W(p_1\cdots p_j-1)
 \bigl(f_x(p_1\cdots p_j-1)-C_x\bigr)=o(x/L).
\end{equation}

Let
\begin{equation}\label{ext:proxy-weight}
 \kappa(m)=\frac1{V(W)\log m}\quad(m\geq2),\qquad
 a(m)=\1_{m>z}\1_{P^-(m)>W}\kappa(m).
\end{equation}
Set $\kappa(1)=0$; values at one are always excluded by the
positive-power threshold tests.
We first replace every prime variable in
\eqref{ext:composite-vanishing} by this coefficient.  Telescope
one variable $m$ at a time and write $n$ for the product of the
other $j-1$ variables.  Dyadically decompose $m,n$ on the support
of $\Phi((mn-1)/x)$.  There are $O(L^2)$ boxes, with
$H_mH_n\asymp x$, and both scales are at least a fixed positive
power of $x$.  The restriction $m>z$ only shortens its dyadic
interval, as permitted in Theorem~\ref{thm:type-II}.
The companion coefficient $\beta_n$ is supported on $W$-rough
integers and is a fixed convolution of prime indicators and proxy
coefficients.  It is bounded by $L^{C_j}\tau_{j-1}(n)$ for a
fixed exponent $C_j$.

The divisor moment bound in Lemma~\ref{lem:moment-bounds} permits
truncation of $\beta_n$ at $L^B$, for a sufficiently large fixed
$B$, with an arbitrarily large log saving in its $\ell^1$ norm.
More explicitly, for any fixed integer $t>1$,
\[
 \sum_{n\asymp H_n}|\beta_n|\1_{|\beta_n|>L^B}
 \leq L^{-B(t-1)}\sum_{n\asymp H_n}|\beta_n|^t
 \ll H_nL^{C_{j,t}-B(t-1)}.
\]
Multiplying by the number and maximum size of the $m$ coefficients
and by the bounded endpoint factor shows that the discarded
terms are $O(xL^{-A})$ for any prescribed $A$, on taking $B$
large enough.  This truncation does not change rough support.

Replacing $\Phi((mn-1)/x)$ by $\Phi(mn/x)$ costs
$O(x^{-1})$ per coefficient, hence only a fixed log power in the
whole sum; it is negligible compared with $x/L$.  Fourier inversion
of the smooth function in $\log(mn/x)$ separates the two variables
into factors $m^{iv}$ and $n^{iv}$, with a rapidly decreasing
Fourier density.  Its tails may be discarded past a fixed log
power, with any desired saving, using the same coefficient bounds.
On the remaining frequencies the discrepancy in the $m$ slot is
exactly the one in Theorem~\ref{thm:type-II}.

To meet its invariance hypothesis globally, clip
$f_x(u)-C_x$ to a fixed interval containing all its values on the
support under consideration, and divide by a fixed bound so that
its modulus is at most one.  Every long-slot prime lies outside
the array bands, so $f_x(pu)=f_x(u)$ for each band prime $p$;
clipping preserves this identity.  Thus the Type II theorem
applies with $F$ equal to this normalized clipped function.
Choose its log-saving target after all the fixed dyadic, Fourier,
and truncation losses.  We then choose $K$ sufficiently large
for that target, simultaneously for all $2\leq j\leq2/\gamma$.
This choice is made after $h,\gamma$, and is independent of the
particular fixed band exponents.  Telescoping all slots gives
\begin{align}\label{ext:proxy-replacement}
 &\sum_{p_1,\ldots,p_j>z}
     \Phi((p_1\cdots p_j-1)/x)
     \mathcal W(p_1\cdots p_j-1)
       (f_x(p_1\cdots p_j-1)-C_x)\notag\\
 &\quad=\sum_{m_1,\ldots,m_j}
     \Phi((m_1\cdots m_j-1)/x)
     \mathcal W(m_1\cdots m_j-1)
       (f_x(m_1\cdots m_j-1)-C_x)
                         \prod_{i=1}^ja(m_i)+o(x/L).
\end{align}

\subsection{Removing candidate prime parts}

Fix $j$ in the preceding finite range.  Construct candidate small
groups from all primes in consecutive bands
$[\exp(y),\exp(2y))$, starting at $y=L^{.28}$ and doubling $y$
while $2y\leq L^{.35}$.  Construct candidate big groups in the
same way, from $y=L^{.40}$ while $2y\leq L^{.45}$.
There are between positive constant multiples of $\log L$
groups of each kind.  They are disjoint, their reciprocal sums
$V_g$ are bounded above and below by positive constants, and they
lie in the allowed small and big ranges of
Theorem~\ref{thm:two-sided}.  The big groups are full prime
intervals.  All candidate primes exceed $W$.

We will group factorizations that differ only in the allocation of
candidate primes among their $j$ factors.  Their coefficients must
agree, so we first remove all candidate prime parts from the threshold
and logarithm in $a(m)$.
For any integer $m$, let $\bar m$ be obtained by removing the
entire prime-power parts belonging to all candidate groups.
Replace $a(m)$ by
\begin{equation}\label{ext:stripped-coefficient}
 \widetilde a(m)=\1_{P^-(m)>W}\1_{\bar m>x^\gamma}
                              \kappa(\bar m).
\end{equation}
We show that this changes the right-hand side of
\eqref{ext:proxy-replacement} by $o(x/L)$ in absolute value.
In particular, the estimate will remain valid before any
cancellation is used.

For these error estimates, decompose tuples into full product
dyads $m_i\asymp H_i$, with $\prod_iH_i\asymp x$ and
$H_i\geq x^\gamma/2$.  There are $O(L^{j-1})$ such boxes.
By Lemma~\ref{lem:rough-counts}, the number of $W$-rough
integers in each dyad is $O(H_iV(W))$.  On an unconditioned
integer dyad,
\begin{align}\label{ext:removed-log-mean}
 \sum_{m\asymp H_i}\log(m/\bar m)
 &\leq\sum_{p\text{ candidate}}\sum_{a\geq1}
                   (\log p)\left\lfloor\frac{2H_i}{p^a}\right\rfloor
 \ll H_i\sum_{p\leq\exp(L^{.45})}\frac{\log p}{p-1}
 \ll H_iL^{.45}.
\end{align}
Thus the number with $\log(m/\bar m)>L^{.9}$ is
$O(H_iL^{-.45})$.  On the support of either coefficient in the
comparison its size is $O((LV(W))^{-1})$.  Taking the bad count
in one coordinate and the rough counts in all others bounds the
bad-tuple mass in a box by
\begin{equation}\label{ext:large-removal-error}
 O\left(\frac{x}{L^j}\frac{L^{-.45}}{V(W)}\right)
 =O(xL^{-j-.21}),
\end{equation}
since $V(W)\asymp L^{-.24}$.  After summing boxes this is
$O(xL^{-1-.21})$.

For the other tuples, a change in the threshold test requires
$\gamma L<\log m_i\leq\gamma L+L^{.9}$ for at least one
coordinate.  There are
$O(L^{.9}L^{j-2})$ product dyads meeting such a strip, because
$j\geq2$.  Their positive mass is $O(x/L^j)$ per box, so the
total is $O(xL^{-1.1})$.  Away from those strips the supports
coincide, and
\[
 \frac{\kappa(\bar m_i)}{\kappa(m_i)}
 =\frac{\log m_i}{\log\bar m_i}=1+O_\gamma(L^{-.1}).
\]
The total positive tuple mass is $O(x/L)$, by the same rough
counts and box count.  Together with the boundedness of
$\mathcal W(f_x-C_x)$ this proves the required absolute
$o(x/L)$ error.

\subsection{Many successes before the signed expansion}

Write $v=m_1\cdots m_j$ and $u=v-1$.  For each candidate group
$g$, independently conditional on the tuple, toss a coin whose
success probability is
\begin{equation}\label{ext:success-probability}
 s_g(u,v)=c_1q_0^{\omega_g(u)-1}\frac{\omega_g(u)}{V_g}
       (q_0/j)^{\omega_g(v)-1}\frac{\omega_g(v)}{V_g}.
\end{equation}
The constant $c_1>0$ is fixed sufficiently small, depending on
$j,q_0$ and the bounds for $V_g$, that these probabilities are at
most one.  The factors with no hit are interpreted as zero.
The division by $j$ in the product-side damping anticipates the $j$
possible allocations of each candidate prime among the factors when
the candidate part of $v$ is squarefree.  Summing those allocations
will recover the $q_0$ damping of the two-sided estimate.
We claim that for some fixed $c_2>0$, outcomes with fewer than
$c_2\log L$ successes in either family have total positive
weighted mass $o(x/L)$ under the coefficient
$\prod_i\widetilde a(m_i)$.

We prove the uniformity needed for this claim on full rough
product dyads.  Select each $m_i$ independently and uniformly
among $W$-rough integers in its full dyad.  For any fixed number
of distinct candidate primes, their product $Q$ satisfies
$\log Q=O(L^{.45})$.  Each prescribed residue class modulo $Q$
is relatively equidistributed among these rough integers, with
error $o(1)$ uniformly in the classes, primes, and dyads.
To verify this, apply Bonferroni to divisibility by primes
$p\leq W$ at depth $D\log L$, where $D$ is a sufficiently
large fixed constant.  The reciprocal sum of those primes is
$O(\log L)$, so taking $D$ large makes the omitted reciprocal
mass smaller than any prescribed log power times $V(W)$.
The moduli in the retained terms are at most
\[
 QW^{O(\log L)}
 =\exp(O(L^{.45}+L^{.24}\log L))=x^{o(1)}.
\]
Since every factor dyad has positive-power length, CRT counting
errors divided by its expected class count tend to zero
uniformly.  Normalizing by the similarly evaluated total rough
count proves the assertion, with arbitrary fixed logarithmic
precision if needed.

In the independent residue model, for a candidate prime $p$ the
events $p\mid u$ and $p\mid v$ are disjoint, with probabilities
\begin{equation}\label{ext:prime-event-probabilities}
 a_p=\frac{(p-1)^{j-1}}{p^j},\qquad
 b_p=1-(1-1/p)^j,
\end{equation}
respectively.  For the first formula all residues must be
nonzero, and the first $j-1$ determine the last.  For the second,
at least one of the $j$ residues must be zero.  Distinct primes
are independent in this model.  In each group,
$\sum_pa_p=V_g+o(1)$ and $\sum_pb_p=jV_g+o(1)$, while the
largest individual probability tends to zero.  Summing the
uniform residue estimates proves convergence of every fixed
joint factorial moment of the two hit counts, for any one or
two groups, uniformly over their choices and over product dyads.
These moments are bounded by $C^r$ at total order $r$, with $C$
depending only on the fixed $j$ and the reciprocal group bounds.

Here is why fixed moments suffice for the bounded probabilities
in \eqref{ext:success-probability}, without any growing-moment
assumption.  First cut each of the at most four counts at a
fixed bound $a$.  The bounded first moments give tail probability
$O(1/a)$ uniformly.  For exact counts below this bound, list the
required hits and impose the absence of further hits by
Bonferroni.  At a fixed depth $b$, the limit superior of the
remainder is at most $C_a C^b/b!$, by the corresponding
factorial moment estimate.  Take $x\to\infty$, then
$b\to\infty$, and then $a\to\infty$.  It follows that the
expectations of the bounded one-group tests and their two-group
products agree with the model up to a uniform $o(1)$.

Explicitly, for a count vector $N=(N_1,\ldots,N_r)$ and a fixed
target $k$, list the $k_i$ required hits of each type.  Conditional
on $N\geq k$ coordinatewise, inclusion--exclusion for no other
hit is the alternating binomial sum in $N_{\rm tot}-|k|$,
multiplied by $\prod_i\binom{N_i}{k_i}$.  The first omitted
term at depth $b$ is bounded by
\[
 \frac{(N_{\rm tot})_{|k|+b+1}}
             {(b+1)!\prod_i k_i!},
\]
where $(n)_r=n(n-1)\cdots(n-r+1)$.  Its expected limit superior
is at most $C^{|k|+b+1}/((b+1)!\prod_i k_i!)$.
This supplies the stated uniform remainder for every exact count
used after the fixed truncation.

In the model, the probability of exactly one hit of each type
in a group is bounded below.  Indeed it equals
\[
 \prod_{p\in g}(1-a_p-b_p)
 \sum_{p\ne q\text{ in }g}
       \frac{a_p}{1-a_p-b_p}\frac{b_q}{1-a_q-b_q},
\]
whose product and sum are bounded below using the displayed
reciprocal bounds and the vanishing largest atom.
On this event the conditional success probability is
$c_1/V_g^2$, also bounded below.  Distinct groups are independent
in the model.  Therefore, if $I_g$ denotes the actual coin
indicator, there are $\beta>0$ and $\epsilon_x\to0$, uniform
in the dyads and groups, such that
\begin{equation}\label{ext:success-covariance}
 \E I_g\geq\beta,\qquad
 |\operatorname{Cov}(I_g,I_{g'})|\leq\epsilon_x\quad(g\ne g').
\end{equation}
Conditional independence of the coins identifies the second
quantity with the covariance of their success probabilities.
For a family of $N\asymp\log L$ groups, Chebyshev gives
\begin{equation}\label{ext:success-chebyshev}
 \Prob\left(\sum_g I_g<\beta N/2\right)
 \leq\frac{4}{\beta^2N}+\frac{4\epsilon_x}{\beta^2}=o(1).
\end{equation}
Only two families occur.  Notice that a uniform $o(1)$ covariance
suffices; no rate relative to $1/\log L$ is required.

Finally transfer this full-dyad probability to the actual
weighted tuples.  A rough product box contains
$O(xV(W)^j)$ tuples, and
$\prod_i\widetilde a(m_i)\ll(LV(W))^{-j}$ on its support.
Thus its failure mass is $o(x/L^j)$ uniformly.  There are
$O(L^{j-1})$ boxes, proving the claimed $o(x/L)$ total.
The bounded factor $\Phi(u/x)\mathcal W(u)(f_x(u)-C_x)$ does not
alter this conclusion.  We discard these failed outcomes now,
while their weights form a nonnegative probability partition.

\subsection{Allocation and the two-sided correlation}

Let $\mathcal C$ be the full candidate set of groups.  The coin
partition is
\[
 1=\sum_{S\subseteq\mathcal C}
       \prod_{g\in S}s_g\prod_{g\notin S}(1-s_g).
\]
Retain only $S$ containing at least $c_2\log L$ groups of each
kind, at the error just proved.  Now expand the failure factors
in each retained term.  Each expanded term has the form
\begin{equation}\label{ext:active-term}
 \pm\prod_{g\in P}s_g,
 \qquad P=S\cup T,\quad T\subseteq\mathcal C\setminus S.
\end{equation}
There are at most $3^{|\mathcal C|}=L^{O(1)}$ terms; their
coefficients, including $c_1^{|P|}$, have fixed log-power bounds.
Every active set $P$ has between fixed positive multiples of
$\log L$ groups of each kind and satisfies all the active-group
hypotheses of Theorem~\ref{thm:two-sided}.  This expansion is
performed after the failed mass has been removed, so it never
multiplies the preceding qualitative $o(x/L)$ error.

For a fixed active set let $s_P=|P|$, let $v_*$ be obtained from
$v$ by removing the entire parts of its active primes, and use
$W_{1,q_0}$ to denote the active one-mark weight with damping
$q_0$.  We may first restrict to $v$ not divisible by the square
of any candidate prime.  Such square exceptions have integer
count
\begin{equation}\label{ext:candidate-square-count}
 O\left(x\sum_{p\text{ candidate}}p^{-2}\right)
 \ll x\exp(-cL^{.28}).
\end{equation}
Their weighted tuple sums remain negligible after all fixed
log-power losses.  Indeed the number of ordered factor tuples
of $v$ is at most $\tau_j(v)$; Cauchy--Schwarz and a fixed
divisor moment bound give an exponential saving in a power of
$L$ for sums of any fixed divisor power over this exceptional
set.  All active marked factors are bounded by fixed log
powers, because there are $O(\log L)$ groups with exponential
damping in each.  The same argument will allow these exceptions
to be restored at the end of the calculation.

Define the core
\begin{equation}\label{ext:allocation-core}
 G_P(v)=\sum_{e_1\cdots e_j=v_*}
       \prod_{i=1}^j
       \left[\1_{P^-(e_i)>W}\1_{\bar e_i>x^\gamma}
                         \kappa(\bar e_i)\right].
\end{equation}
Here the bar still removes \emph{all} candidate prime parts,
including the inactive ones.  This definition implies
$G_P(v)=G_P(v_*)$ and
\begin{equation}\label{ext:core-bound}
 |G_P(v)|\leq (C/(LV(W)))^j\tau_j(v_*)
               \leq L^{C_j}\tau(v_*)^{C_j}.
\end{equation}
The second bound follows from the elementary inequality
$\tau_j(n)\leq\tau(n)^{j-1}$, obtained prime by prime.

Off the square exceptions, fix a factorization
$e_1\cdots e_j=v_*$.  Every active prime dividing $v$ can be
assigned independently to any of the $j$ slots.  These are all
the original factorizations with that stripped factorization,
and all have the same coefficient in
\eqref{ext:stripped-coefficient}: active primes exceed $W$ and
the tests and logarithms depend only on the bars.  Thus there
are $j^{\omega_P(v)}$ equal contributions.  The exact identity
\begin{equation}\label{ext:allocation-identity}
 j^{\omega_P(v)}(q_0/j)^{\omega_P(v)-s_P}
     =j^{s_P}q_0^{\omega_P(v)-s_P}
\end{equation}
shows that summing the $v$-side factor in
\eqref{ext:active-term} over factorizations gives
\[
 j^{s_P}W_{1,q_0}(v)G_P(v).
\]
This is an identity of weights, including the mark factors
$\prod_g\omega_g(v)/V_g$.

Restoring the square exceptions by
\eqref{ext:candidate-square-count} and its divisor-moment bound,
each expanded term is, up to a negligible error, a constant of
fixed log-power size times
\begin{equation}\label{ext:correlation-to-use}
 \sum_{u\geq1}\Phi(u/x)\mathcal W(u)(f_x(u)-C_x)
        W_{1,q_0}(u)W_{1,q_0}(u+1)G_P(u+1).
\end{equation}
Take $\Psi(y)=y\Phi(y)$.  Since
$\Phi(u/x)=x\Psi(u/x)/u$, Theorem~\ref{thm:two-sided} bounds
\eqref{ext:correlation-to-use} by $O_A(xL^{-A})$ for every
fixed $A$.  Its invariant core hypothesis is exactly
\eqref{ext:allocation-core}, and its growth hypothesis is
\eqref{ext:core-bound}.  Its other endpoint is exactly
$F_x(u)=\mathcal W(u)(f_x(u)-C_x)$.  Choose $A$ after the fixed
costs from $3^{|\mathcal C|}$, $j^{s_P}$, and the coefficients.
They are all absorbed.  Summing the expanded terms and adding
back the earlier, unamplified, discarded masses proves that the
proxy sum in \eqref{ext:proxy-replacement} is $o(x/L)$.
This proves \eqref{ext:composite-vanishing}.

\subsection{Removal of the small-band marks}

Subtract the composite contribution from
\eqref{ext:presift}.  All primes in the support exceed $z$ for
large $x$, and we have proved
\begin{equation}\label{ext:weighted-primes}
 \limsup_{x\to\infty}\frac Lx
 \left|\sum_{p\text{ prime}}\Phi((p-1)/x)
       \frac{\mathcal W(p-1)}{m_{\mathcal W}}
                    (f_x(p-1)-C_x)\right|
 \ll e^{-h}/\gamma.
\end{equation}
The constant here remains independent of $K,h$.

Fix this $h$ and its sufficient fixed $K$.  For any finite
$b\geq1$ choose $b$ disjoint arrays, each with $K$ distinct
exponents, using distinct exponents across all arrays inside
$(.1,.2)$.  This is possible for every finite $b$, and the
Type II choice of $K$ applies to every one of them.  Write
$\mathcal W_a,m_a$ for their weights and model means, and set
\begin{equation}\label{ext:deweight-square}
 Z_b(u)=\left(\frac1b\sum_{a=1}^b
                     \frac{\mathcal W_a(u)}{m_a}-1\right)^2.
\end{equation}
Disjoint arrays are independent in the divisor model.  Their
normalized weights have mean one and uniformly bounded second
moments by \eqref{ext:model-bounds}.  Hence
\begin{equation}\label{ext:deweight-model}
 M_x(Z_b)=\frac1{b^2}\sum_{a=1}^b
               \operatorname{Var}_M(\mathcal W_a/m_a)
           \leq C_K/b.
\end{equation}
The square and all its cross terms lie within
Lemma~\ref{lem:progressions}, whose coefficient bounds are
uniform for these normalized weights when $b,K$ are fixed.

Apply the nonnegative upper bound of the block sieve at the
fixed depth $h'=2$, with $z'=x^{c_0/20}$, to the weight
$\Phi(u/x)Z_b(u)$.  Its remainder range is
\[
 (z')^{4h'+2}=x^{c_0/2}<x^{c_0},
\]
so Lemma~\ref{lem:progressions} without tuple factors applies.
Every prime $u+1$ in the support exceeds $z'$.  The sieve main
term is bounded by a fixed constant times
$xM_x(Z_b)/\log z'$, and its remainder is $o(x/L)$.
It follows that
\begin{equation}\label{ext:prime-deweight-square}
 \limsup_{x\to\infty}\frac Lx
     \sum_{p\text{ prime}}\Phi((p-1)/x)Z_b(p-1)
     \ll C_K/b.
\end{equation}
The implied constant here depends on the fixed $c_0$ and cutoff,
but not on $b,K$; the latter dependence is recorded in $C_K$.

Average \eqref{ext:weighted-primes} over the $b$ arrays.  By
Cauchy--Schwarz, the difference between this averaged weighted
sum and the unweighted sum is at most
\begin{align*}
 &\left(\sum_p\Phi((p-1)/x)Z_b(p-1)\right)^{1/2}\\
 &\hspace{12mm}\cdot
 \left(\sum_p\Phi((p-1)/x)(f_x(p-1)-C_x)^2\right)^{1/2}.
\end{align*}
The second factor is $O((x/L)^{1/2})$ by the boundedness of
$f_x-C_x$ and the prime number theorem.  We conclude that
\begin{equation}\label{ext:final-limsup}
 \limsup_{x\to\infty}\frac Lx
 \left|\sum_p\Phi((p-1)/x)(f_x(p-1)-C_x)\right|
 \leq C\frac{e^{-h}}\gamma+C\sqrt{C_K/b}.
\end{equation}
All constants in the preceding remainder estimates were allowed
to depend on the finite $b$; they have disappeared in this real
$x$ limit.  First let $b$ be arbitrarily large with $h,K$ fixed.
Then let the even integer $h$ be arbitrarily large, choosing its
new sufficient fixed $K$ each time.  Since
$\gamma^{-1}=(4h+3)/c_0$, the surviving bound tends to zero.
This proves Theorem~\ref{thm:interior}.

In particular, the order of choices is: fix the interior data
and $c_0$; fix $h$, then $\gamma$ and the finitely many composite
lengths; choose the Type II precision and $K$; fix finitely many
disjoint exponent arrays; take the limit in real $x$; then
remove the arrays by $b\to\infty$ and the sieve error by
$h\to\infty$.  No number of arrays grows with $x$ in an
application of either correlation theorem or the progression
lemma.
